\documentclass[11pt]{article}
\usepackage{ap-lang}
\title{Part 1: The Rhetorical Situation: Question Set}
\author{Paul Wang}
\date{November 2026}
\docpart{Part 1 \textperiodcentered\ Question Set}

\begin{document}
\maketitle

\begin{abstract}
Three exam-style passages with multiple-choice questions, followed by a full rhetorical analysis essay (Question 2). Passages 1 and 2 are reading sets: a speech to a county board and a letter to a fishing cooperative, both written to audiences that may resist the writer's request. Passage 3 is a writing set built on a student's draft letter with numbered sentences. Most questions test the rhetorical situation (Part 1); Rule tags beginning 2, 3, or 4 point to the guide for that part. All passages, writers, places, and figures are fictional and were written for this practice set. Allow about 35 minutes for the multiple choice and 40 minutes for the essay.
\end{abstract}

\scoringline{MCQ: \blank\ / \qtotal \qquad Essay: \blank\ / 6 \qquad Target: 20+ MCQ and 4+ essay \qquad Pace: 80 s per MCQ}

\section{Passage 1: The bookmobile (reading)}

\begin{passage}{Passage 1}{The following passage is from remarks delivered in 2019 by Margaret Odell, a retired elementary school teacher, at a public meeting of a rural county's Board of Supervisors. The board was considering a budget that would eliminate the county library's bookmobile.}
\pn{1} Good evening. My name is Margaret Odell, and for thirty-one years I taught fourth grade at Pine Hollow Elementary, which means that a good number of you on this board once sat in a small chair in my classroom and told me, with a straight face, that you had finished your reading. I forgave you then, and I am prepared to forgive you now, provided you vote correctly tonight.

\pn{2} I am here about the bookmobile. I have read the budget memo, all eleven pages of it, and I want to say first that I understand it. The county has lost two employers in three years. The road fund is short. Every line in that memo was written by someone trying to make an impossible sum come out even, and I do not envy you the arithmetic.

\pn{3} But I would ask you to consider what the bookmobile line actually buys. It does not buy a truck. It buys Tuesday afternoons at the Fenwick crossroads, where eleven children wait by the feed store because their nearest library is forty minutes away and their parents work until six. It buys the large-print mysteries that Mr.\ Albert Dunn, who is ninety-one and no longer drives, reads at a rate of three a week. It buys the only place in the eastern half of this county where a person can print a job application without asking a favor.

\pn{4} The memo calls the bookmobile a ``low-utilization service.'' I suppose that is true if you measure it by the number of people who use it. Measure it instead by the number of people who have nothing else, and it is the most utilized thing this county owns.

\pn{5} You will hear that the families out east can drive to the main branch. Some can. The ones who most need it cannot, and a cut that falls only on the people least able to absorb it is not a saving. It is a transfer of the burden from this table to theirs.

\pn{6} So I am asking you to keep the bookmobile for one more year, and I am offering something in return. The Friends of the Library will raise the cost of its fuel for that year ourselves. We have already begun. I sold my first pie this afternoon, to your county clerk, at a price I will not repeat in public.
\end{passage}

\begin{q}{3}{1.5 Accommodating audience}{She uses her personal history with the board to set a warm, familiar tone before her request.}
In paragraph 1, Odell's reference to board members as her former students primarily serves to
\ch{suggest that the board members cannot be trusted to read the budget memo}{establish her expertise in public finance}{create a familiar, good-humored relationship with an audience she knows personally}{signal that her remarks will not concern the budget}
\end{q}

\begin{q}{2}{1.1 Exigence}{The board's vote on a budget cutting the bookmobile is the event that prompts the speech.}
The exigence for Odell's remarks is
\ch{Odell's retirement after thirty-one years of teaching}{the board's consideration of a budget that would eliminate the bookmobile}{a decline in reading among the county's children}{the county's loss of two employers in three years}
\end{q}

\begin{q}{2}{1.3 Ethos}{Granting the difficulty of the board's task makes her request seem reasonable and fair.}
In paragraph 2, Odell's statement that she understands the budget memo mainly serves to
\ch{agree that the bookmobile should be eliminated}{acknowledge the board's difficult position so that her request appears fair-minded}{criticize the staff who wrote the memo}{demonstrate her training in accounting}
\end{q}

\begin{q}{1}{4.2 Syntax}{Repeated ``It buys'' turns an abstract budget line into concrete services for specific people.}
The repetition of ``It buys'' in paragraph 3 primarily serves to
\ch{redefine a budget line as specific services to specific residents}{list the mechanical features of the bookmobile}{shift the passage to a hostile tone}{emphasize how expensive the bookmobile is to operate}
\end{q}

\begin{q}{1}{2.4 Counterargument}{She accepts the phrase's literal accuracy, then changes the measure of value.}
In paragraph 4, Odell responds to the memo's description of the bookmobile as a ``low-utilization service'' by
\ch{granting its literal accuracy while proposing a different measure of the service's value}{proposing a new formula for calculating the county budget}{rejecting the memo's figures as inaccurate}{mocking the author of the memo}
\end{q}

\begin{q}{4}{2.1 Claims}{The cut shifts costs onto residents least able to bear them rather than removing them.}
The claim in paragraph 5 that the cut ``is a transfer of the burden from this table to theirs'' suggests that eliminating the bookmobile would
\ch{require board members to pay for the service personally}{move the bookmobile's services to the main branch}{save the county money over the long term}{shift costs onto vulnerable residents rather than eliminate them}
\end{q}

\begin{q}{3}{2.2 Evidence}{Named people and places give human faces to the usage figures.}
The details about the children at the Fenwick crossroads and Mr.\ Albert Dunn (paragraph 3) function mainly to
\ch{suggest that the bookmobile primarily serves elderly residents}{prove that the bookmobile's usage numbers are higher than the memo claims}{give concrete, human examples of what the cut would take away}{provide expert testimony about rural libraries}
\end{q}

\begin{q}{2}{1.5 Accommodating audience}{Sharing the cost makes the request easier to grant, and the joke echoes the warm opening.}
In paragraph 6, Odell's offer and her remark about the pie mainly serve to
\ch{suggest that the county clerk acted improperly}{make her request easier to grant while ending on the good humor of her opening}{show that private fundraising can permanently replace county funding}{change the subject from the bookmobile to fundraising}
\end{q}

\begin{q}{4}{4.4 Tone}{Warm humor throughout, but she presses a clear demand.}
Odell's tone throughout the passage is best described as
\chh{detached and clinical}{bitter and accusatory}{pleading and desperate}{affectionate but firm}
\end{q}

\section{Passage 2: The Ledge (reading)}

\begin{passage}{Passage 2}{The following passage is from a letter published in 2023 in the newsletter of a fishermen's cooperative in a small coastal town. The writer, a marine biologist who grew up in the town, is responding to members' anger over a proposed two-year closure of the cod fishery on the Ledge, a fishing ground offshore.}
\pn{1} I know what some of you said at the co-op meeting, because my uncle called me afterward to repeat it, word for word, with what I can only describe as enthusiasm. I will not argue with the parts about my character. I will argue with the rest.

\pn{2} I was born four streets from the harbor. My grandfather fished the Ledge for forty years, and the first thing I ever counted was cod on the deck of his boat. I mention this not to claim I am one of you---you would correct me, and you would be right---but so you know that when I say the closure is necessary, I am not saying it from a distance.

\pn{3} Here is what the surveys show. In 1990 a two-hour tow on the Ledge brought up, on average, 310 pounds of cod. Last spring the same tow brought up 41. The fish we are catching are also younger: most have not yet spawned even once. A fishery can survive hard years. It cannot survive taking the fish before they replace themselves, because then every hard year is followed by a harder one.

\pn{4} You have heard this before. You heard it in 1994, when the scientists were right, and in 2008, when, I will admit, we were wrong about how quickly the haddock would come back. I am not asking you to trust scientists because we are always right. I am asking you to look at your own logbooks. They tell the same story as ours, only in your handwriting.

\pn{5} A closure is not a punishment. It is a loan from the present to the future, and like any loan, it is hardest to make when money is short. That is why the proposal includes payments to boats that stay in port, and why I have asked the state to put a fishermen's representative on the board that decides when the Ledge reopens. You should not have to take our word for when it is ready. You should be there to see it.
\end{passage}

\begin{q}{3}{1.2 Audience}{Light humor acknowledges the audience's anger and lowers the temperature.}
In paragraph 1, the writer's comment that her uncle repeated the meeting's remarks ``with what I can only describe as enthusiasm'' primarily serves to
\ch{dismiss the members' concerns as trivial}{reveal a family conflict that becomes the letter's subject}{acknowledge the members' anger with humor that defuses hostility}{establish her scientific authority}
\end{q}

\begin{q}{1}{1.3 Ethos}{Personal roots build credibility without claiming insider status.}
Paragraph 2 primarily serves to
\ch{establish the writer's connection to the town without overstating it}{claim that she is a working fisherman}{describe her grandfather's fishing methods}{argue that the writer's family history proves the closure is necessary}
\end{q}

\begin{q}{3}{1.5 Accommodating audience}{She concedes the members' sense of who belongs, showing respect.}
The phrase ``you would correct me, and you would be right'' (paragraph 2) suggests that the writer
\ch{believes the members cannot be persuaded by an outsider}{wants to mock the members' pride}{respects the members' sense of who counts as one of them}{doubts the accuracy of her own conclusions}
\end{q}

\begin{q}{4}{2.2 Evidence}{Survey figures (310 vs.\ 41 pounds) document the decline that justifies the closure.}
The figures in paragraph 3 function primarily to
\ch{suggest that fishing technology has improved since 1990}{show that earlier scientific predictions failed}{prove that haddock recover quickly after a closure}{provide concrete evidence of the decline that makes the closure necessary}
\end{q}

\begin{q}{2}{3.2 Methods of development}{The paired sentences contrast hard years with overfishing to isolate the cause of collapse.}
The sentences ``A fishery can survive hard years. It cannot survive taking the fish before they replace themselves'' (paragraph 3) primarily
\ch{suggest that the fishery has already collapsed}{distinguish ordinary hardship from the specific cause of lasting decline}{shift the letter from evidence to personal anecdote}{introduce a counterargument the writer then rejects}
\end{q}

\begin{q}{3}{2.4 Concession}{Admitting a past error to readers who remember it strengthens credibility.}
The writer's admission in paragraph 4 that scientists ``were wrong'' in 2008 most likely
\ch{undermines her argument for the closure}{accuses the members of ignoring science}{strengthens her credibility with readers who remember the error}{changes the subject from cod to haddock}
\end{q}

\begin{q}{1}{2.2 Evidence}{The members' own records confirm the surveys, so the conclusion is theirs too.}
The statement that the logbooks ``tell the same story as ours, only in your handwriting'' (paragraph 4) suggests that
\ch{the members' own records confirm what the surveys show}{the members' logbooks are less reliable than scientific surveys}{the scientists copied their data from the logbooks}{the members have kept inaccurate records}
\end{q}

\begin{q}{2}{4.3 Comparisons}{The loan analogy reframes closure as a temporary sacrifice with a future return.}
The comparison of the closure to ``a loan from the present to the future'' (paragraph 5) primarily serves to
\ch{advise members to borrow money during the closure}{reframe the closure as a temporary sacrifice that will be repaid}{emphasize that the closure's costs outweigh its benefits}{show that the state will pay for every lost catch}
\end{q}

\begin{q}{1}{1.5 Accommodating audience}{A seat on the reopening board answers the members' distrust of scientists.}
The final two sentences of the letter mainly serve to
\ch{offer the members a role in deciding the outcome, addressing their distrust}{suggest that the scientists cannot be trusted}{concede that the closure may become permanent}{warn the members that they will be punished if they resist}
\end{q}

\section{Passage 3: Saturday library hours (writing)}

\begin{passage}{Passage 3}{The following is a draft of a letter that a student is writing to her school's principal. It may need revision.}
\pn{1} \sn{1} Every Saturday morning, a line forms outside the Westbrook Public Library before it opens. \sn{2} Many of the people in that line are students from our school. \sn{3} They come for quiet, for Wi-Fi, and for tables big enough to spread out a chemistry lab report. \sn{4} Since the public library cut its weekend hours in January, however, it has been open only from ten to noon, and students are now turned away once every seat is taken.

\pn{2} \sn{5} I am writing to propose that our school open its library from 9 a.m.\ to 1 p.m.\ on Saturdays. \sn{6} Our library already has forty-two seats, twenty computers, and a printer that students can use for free. \sn{7} For most of the weekend, all of it sits behind a locked door.

\pn{3} \sn{8} The main cost would be staffing. \sn{9} The library would need one adult supervisor for four hours each week. \sn{10} Parent volunteers from the PTA have already offered to share these shifts, and the student council has voted to help shelve books and check out laptops. \sn{11} The custodial staff already works on Saturdays to prepare the building for weekend sports events.

\pn{4} \sn{12} Some teachers have worried that students will use the space to socialize rather than to study. \sn{13} That is a fair concern. \sn{14} A trial run would answer it: if attendance or behavior is poor after eight weeks, the program could simply end.

\pn{5} \sn{15} Students at our school want to do their work. \sn{16} We are asking for a place to do it.
\end{passage}

\begin{q}{2}{2.2 Evidence}{A specific student's experience shows the effect of the reduced hours.}
The writer wants to add a sentence after sentence 4 to emphasize how the reduced hours affect students. Which choice best accomplishes this goal?
\ch{Many people enjoy reading on the weekend.}{Last month, a junior told me she waited forty minutes in the cold, was turned away, and finished her essay on her phone in the parking lot.}{Public libraries across the country face budget pressures.}{The Westbrook Public Library was built in 1962 and renovated in 2005.}
\end{q}

\begin{q}{1}{4.1 Diction}{``Sits behind a locked door'' sharpens the contrast between rich resources and no access.}
In sentence 7 (reproduced below), the writer wants to emphasize the contrast between the library's resources and their lack of use on weekends.

\emph{For most of the weekend, all of it sits behind a locked door.}

Which version of sentence 7 best accomplishes this goal?
\ch{(as it is now)}{The library is closed on weekends.}{On weekends, the library is not open for students to use.}{Some weekend days, the library may be closed.}
\end{q}

\begin{q}{4}{2.2 Evidence}{Sentence 11 removes a cost objection: the building is already open and staffed.}
The writer is considering deleting sentence 11. Should the writer keep or delete this sentence?
\ch{Delete it, because custodial staff have no connection to the library.}{Delete it, because it introduces a new topic, weekend sports.}{Keep it, because it gives the main reason students need the library.}{Keep it, because it shows the building would already be open, which answers a likely concern about cost.}
\end{q}

\begin{q}{4}{3.4 Transitions}{The trial run resolves the concern just granted: ``Fortunately'' signals a welcome answer.}
Which transition, if placed at the beginning of sentence 14 (adjusting capitalization as needed), best connects it to sentence 13?
\chh{As a result,}{For instance,}{Similarly,}{Fortunately,}
\end{q}

\begin{q}{3}{1.5 Accommodating audience}{Naming existing resources answers a principal's likely concern about cost.}
The writer wants sentence 5 to state her proposal in a way that anticipates the principal's likely concern about cost. Which version of sentence 5 best accomplishes this goal?
\ch{(as it is now)}{I am writing to propose that our school open its library from 9 a.m.\ to 1 p.m.\ on Saturdays, because I think it would be nice.}{I am writing to propose that our school open its library from 9 a.m.\ to 1 p.m.\ on Saturdays, a change that would put resources we already have to use.}{I am writing to propose that our school open its library from 9 a.m.\ to 1 p.m.\ on Saturdays, even though many teachers oppose the idea.}
\end{q}

\begin{q}{1}{2.2 Evidence}{Specific signed-up volunteers show the staffing plan is real.}
The writer wants to add evidence to paragraph 3 that the staffing plan is realistic. Which choice best accomplishes this goal?
\ch{Twelve parents have already signed up for shifts through March, enough to cover every Saturday.}{Many schools in other states have libraries.}{The PTA held a successful bake sale last fall.}{Students say the library's chairs are comfortable.}
\end{q}

\begin{q}{4}{3.3 Structure}{Ending on a reversible trial reminds the principal how little the proposal risks.}
The writer wants to end the letter by reminding the principal how little risk the proposal carries. Which choice, used to replace sentence 16, best accomplishes this goal?
\ch{(as it is now)}{The public library should restore its weekend hours.}{Libraries have existed since ancient times.}{Give us eight Saturdays, and if we do not use them well, you can lock the door again.}
\end{q}

\begin{q}{2}{3.3 Structure}{A cost figure belongs right after the staffing requirement it prices.}
The writer wants to add the following sentence to the letter.

\emph{That would cost about ninety dollars a week if the school paid a staff member instead.}

Where would the sentence best be placed?
\chh{After sentence 6}{After sentence 9}{After sentence 12}{After sentence 15}
\end{q}

\section{Free response}

\begin{frq}{Question 2 \textperiodcentered\ Rhetorical Analysis \textperiodcentered\ 6 points \textperiodcentered\ 40 minutes}
\small
\textbf{Directions.} The following speech was delivered in 1977 by Dorothy Kaye, the water commissioner of Alder Creek, a small city in California, at a public meeting during the second year of a severe drought. (Kaye and Alder Creek are fictional.) Read the passage carefully.

\textbf{Prompt.} Write an essay that analyzes the rhetorical choices Kaye makes to persuade residents to reduce their water use voluntarily.

In your response you should do the following:
\begin{itemize}[leftmargin=1.5em,itemsep=1pt]
\item Respond to the prompt with a thesis that analyzes the writer's rhetorical choices.
\item Select and use evidence to support your line of reasoning.
\item Explain how the evidence supports your line of reasoning.
\item Demonstrate an understanding of the rhetorical situation.
\item Use appropriate grammar and punctuation in communicating your argument.
\end{itemize}
\end{frq}

\begin{passage}{Question 2 passage}{Dorothy Kaye, remarks at a public meeting, Alder Creek, California, 1977.}
\pn{1} Thank you for coming out on a Tuesday, and thank you to whoever brought the lemonade. I noticed it was made with very little water. That is the kind of civic spirit I am here to talk about.

\pn{2} Our reservoir is at nineteen percent. I want you to hold that number in your head the way you would hold the balance of your checking account in the last week of the month. Nineteen percent is not empty. It is also not a number anyone in this room would be comfortable seeing on a bank statement with the rent due.

\pn{3} The council could order rationing tonight. There are cities up the valley that have done exactly that, with meters and fines and inspectors walking the alleys at dawn to look for wet sidewalks. I do not want to send an inspector to your house. I suspect you do not want to meet him. And I do not believe this city needs one.

\pn{4} I believe this because I have seen what this town does when the creek runs low. In 1924, when the wells failed, families on the east side carried water to the west side in milk cans until the rains came. My mother was one of the girls carrying those cans. She told me that nobody ordered them to do it. It simply did not occur to anyone that a neighbor's thirst was someone else's problem.

\pn{5} What I am asking of you is smaller than that. Turn off the tap while you brush your teeth. Let the lawn go gold; it will come back, and gold is a perfectly respectable color for California in August. Run the dishwasher when it is full and not before. Keep a bucket in the shower. None of this is heroic. All of it, done in ten thousand houses, is the difference between a hard summer and a dangerous one.

\pn{6} Every Sunday I will print the reservoir level in this paper, along with our use, street by street. Not to shame anyone. To let you see what you are doing. When a street cuts its use by a quarter, I will say so, by name.

\pn{7} The rain will come back. It always has. The question is only what kind of town we will have been while we waited for it.
\end{passage}

\makekey

\subsection*{Essay scoring guide (6 points)}
\small
The rows follow the published AP English Language rubric (thesis 1, evidence and commentary 4, sophistication 1). The ``earns credit'' column shows what a top response does with this passage.
\begin{compare}{Row}{Earns credit}{Common miss}
A. Thesis (1) & A defensible thesis that analyzes Kaye's choices in relation to her purpose, e.g.\ ``Kaye persuades residents to conserve by translating the crisis into the familiar terms of a household budget, holding up the town's own history of generosity as the standard to meet, and keeping her requests so small and specific that refusing them would feel unneighborly.'' & ``Kaye uses ethos, pathos, and logos to persuade the town''; restating the prompt; summarizing the speech. \\
B. Evidence and commentary (4) & Specific quotations support every claim, and commentary consistently explains how each choice works on this audience: the checking-account comparison (paragraph 2) makes an abstract percentage personal and urgent without panic; the contrast with inspectors ``at dawn'' (paragraph 3) makes voluntary conservation the attractive alternative; the 1924 milk cans and her mother's story (paragraph 4) define conservation as the town's character, so cooperating means being true to themselves; the list of small imperatives and ``none of this is heroic'' (paragraph 5) make compliance feel easy; public street-by-street reporting (paragraph 6) adds friendly accountability. 4 points: all claims supported and consistently explained. 3: explanations for only some. 2: some evidence, little line of reasoning. 1: summary. & Listing devices with a quotation each and no explanation of effect; paraphrasing the speech paragraph by paragraph; claiming she ``threatens'' residents (she explicitly rejects inspectors). \\
C. Sophistication (1) & Explains a tension Kaye manages, e.g.\ she holds the threat of rationing in reserve while insisting she does not want to use it, so the warning works as an incentive rather than a command; or shows how the final sentence reframes conservation as a question of the town's identity, turning a utility request into civic memory; or writes with a consistently vivid, controlled style. & One clever sentence in the conclusion; a broad claim about all droughts or all politicians with no link to the passage. \\
\end{compare}

\end{document}
